\section{Unbounded stacks: productive summaries and closed exclusions}
\label{sec:pushdown}

\status{Implemented and tested in Python, including the regular-target compiler.
The suffix bridge, summary checker soundness, and source reifier have not been
implemented or checked in Lean.}

\subsection{Input language and the question being decided}
A normalized pushdown system is a tuple $(Q,\Gamma,\Delta)$ with finite control
set $Q$, finite stack alphabet $\Gamma$, and rules
\[
   (p,A)\longrightarrow(q,w),\qquad |w|\leq 2.
\]
The top of the stack is its \emph{first} symbol. A rule induces
\[
   (p,A\beta)\longrightarrow(q,w\beta)
\]
for every suffix $\beta$. There is no transition from an empty stack. The input
also specifies an initial configuration $(p_0,\alpha)$ and a set of final
controls $F$. The basic query is
\begin{equation}
   \exists f\in F:\quad (p_0,\alpha)\longrightarrow^*(f,\eps).
   \label{eq:pds-query}
\end{equation}
Finite control and alphabet do not make the configuration space finite. Rules
may repeatedly increase the stack height. The worker does not impose a height,
run-length, or recursion-depth bound.

The length-two restriction is part of the accepted input language, not an
unchecked preprocessing convention. A future source adapter can normalize
other transition syntax, but it must prove its own correspondence theorem.
Arbitrary data-bearing frames, two stacks, higher-order stacks, and unbounded
integer control variables are not accepted simply because a program happens
to use recursion.

This is a specialization of the established saturation approach to pushdown
reachability \cite{carayol-hague}. The specialization makes both the certificate
and its source theorem especially small.

\subsection{The finite object that summarizes an infinite configuration space}
Write $S(p,A,q)$ when the one-symbol stack can be removed:
\[
   S(p,A,q)\quad\Longleftrightarrow\quad
   (p,A)\longrightarrow^*(q,\eps).
\]
There are only $|Q|^2|\Gamma|$ possible triples. The essential feature of this
relation is that it describes a transformation valid above \emph{every}
suffix, rather than a finite selection of suffixes.

\begin{lemma}[Suffix lifting]
\label{lem:suffix}
If $S(p,A,q)$, then for every $\beta$ there is a run
$(p,A\beta)\longrightarrow^*(q,\beta)$ that does not inspect or modify
$\beta$ before its final configuration.
\end{lemma}
\begin{proof}
Take a run from the one-symbol configuration to the empty stack. Every
configuration before its final one has a nonempty stack: after becoming empty
it could make no further step. Append $\beta$ to every stack in the run. Each
step remains an instance of the same rule, because rules inspect only the top
symbol and preserve the rest. The appended suffix is therefore untouched.
\end{proof}

\begin{remark}[Do not weaken the definition]
An arbitrary run from $(p,A\beta)$ to $(q,\beta)$ might first expose, inspect,
and later reconstruct $\beta$. Such a run is not by itself a suffix-parametric
pop summary. The definition uses the one-symbol run, and the bridge proves
protected suffix lifting. This distinction matters when translating recursive
programs with stack-dependent behavior.
\end{remark}

The summary relation is generated by three local rules:
\begin{align}
\frac{(p,A)\to(q,\eps)\in\Delta}{S(p,A,q)}
&\qquad\text{(pop)},\label{eq:pop}\\[0.2em]
\frac{(p,A)\to(r,B)\in\Delta\quad S(r,B,q)}{S(p,A,q)}
&\qquad\text{(one child)},\label{eq:unary}\\[0.2em]
\frac{(p,A)\to(r,BC)\in\Delta\quad S(r,B,m)\quad S(m,C,q)}{S(p,A,q)}
&\qquad\text{(two children)}.\label{eq:binary}
\end{align}
These rules describe a least fixed point. A relation containing a circular
justification is not automatically evidence of reachability.

\begin{theorem}[Exactness of pop-summary closure]
\label{thm:pop-exact}
The least relation closed under \eqref{eq:pop}--\eqref{eq:binary} is precisely
$S$. Consequently \eqref{eq:pds-query}, for
$\alpha=A_1\cdots A_\ell$, holds exactly when there are controls
\[
 q_0=p_0,\qquad q_\ell\in F,\qquad q_i\in Q\ (0\leq i\leq\ell)
 \quad\text{with}\quad S(q_{i-1},A_i,q_i)\quad(1\leq i\leq\ell),
\]
Only the last control is required to belong to $F$.
\end{theorem}
\begin{proof}
For soundness, induct on a finite derivation. A pop is a real transition. A
one-child rule is followed by the child's run. For a two-child rule, lift the
first child's run over the suffix $C$, then run the second child. Thus every
derived triple describes a genuine pop.

For completeness, induct on the number of transitions in a run removing a
one-symbol stack. Inspect its first rule. A zero-symbol replacement finishes
the run, because the empty stack cannot move. A one-symbol replacement leaves
a strictly shorter pop run. A two-symbol replacement produces $BC$. Before
$C$ can be removed, the part above it must disappear; take the first moment
when $C$ is exposed. Up to this moment the machine cannot have inspected $C$.
Removing this protected suffix gives a run that pops $B$ to some control $m$.
The remaining segment pops $C$. Both segments are shorter than the original
run, so induction supplies both premises of \eqref{eq:binary}.

For a longer initial word, decompose a run at the first exposure of each
successive original suffix. This gives the controls $q_i$. Conversely, compose
the summary runs using Lemma~\ref{lem:suffix}. For $\ell=0$, the criterion is
just $p_0\in F$, correctly accepting a zero-step run.
\end{proof}

\subsection{Search: indexed semi-naive saturation}
The executable solver maintains a table of discovered triples and an agenda of
new triples. A triple is inserted once. Alongside it, the solver records one
rule identifier and the zero, one, or two earlier triples used to derive it.

Three indexes avoid rescanning the whole rule set. A unary rule is indexed by
its required child head $(r,B)$. A binary rule is indexed both by its first
child head $(r,B)$ and by its second stack symbol $C$. A fourth index maps a
summary head $(m,C)$ to all currently known endpoints. When a new summary
arrives, it can therefore act as either child of a binary rule. Both directions
are necessary: the first child may be found after the second.

\begin{lstlisting}[caption={Agenda skeleton; the accompanying implementation includes both join directions.}]
insert all pop-rule summaries, with zero children
while agenda is nonempty:
    take a fresh summary (r, B, m)
    for every unary rule requiring (r, B):
        insert its conclusion using this child
    for every binary rule requiring (r, B) as its first child:
        join with all known (m, C, q)
        insert conclusions with the two child identifiers
    for every binary rule whose second symbol is B:
        look up its first child ending at r
        when present, join it with the new (r, B, m)
query the completed relation along the exact initial stack
return a sliced productive DAG, or a closed exclusion relation
\end{lstlisting}

Let $n=|Q|$, $a=|\Gamma|$, $r_i$ be the number of rules replacing the top
symbol by $i$ symbols, and $M\leq n^2a$ the number of summaries. The core
indexed closure performs
\[
 O(M+r_0+r_1n+r_2n^2)
\]
insertions, joins, and table operations under constant-time lookup assumptions.
The bound counts combinations of a rule, a middle control, and an endpoint;
there are no unbounded configuration objects. Storage for the closure and
provenance is $O(M+|\Delta|)$ records. Deterministic sorting, dictionary
implementation costs, and Python object overhead are additional.

Querying an initial stack of length $\ell$ is a finite relational composition,
with $O(\ell n^2)$ endpoint tests in a simple representation. The prototype
copies short root paths while doing this query; a production implementation
should retain predecessor links instead, eliminating quadratic path-copying
costs on long initial words. These implementation details do not affect the
fixed-point termination proof.

\subsection{Positive certificate: productivity, not just closure}
A positive certificate contains a topologically ordered array of nodes. Node
$i$ records a triple $(p,A,q)$, an input rule identifier, and child identifiers
strictly smaller than $i$. The checker obtains the rule from the separately
supplied problem, never from a rule description inside the certificate.
It checks that the rule's source matches $(p,A)$, that its replacement symbols
match the children in order, and that the children's control endpoints form a
chain ending at $q$.

A separate root list has exactly one node identifier per symbol in the actual
initial stack. It must start at the actual initial control and end in the actual
final set. The checker rejects duplicate summary triples, forward references,
self-references, illegal identifiers, omitted root symbols, and mismatched
control chains. It does not accept a certificate-supplied claim that the target
has been reached.

\begin{proposition}[Positive replay soundness]
If this checker accepts a positive certificate for a valid input, then
\eqref{eq:pds-query} holds.
\end{proposition}
\begin{proof}
Induct on the node index. Every child has already been justified, and the local
rule check is exactly one of \eqref{eq:pop}--\eqref{eq:binary}. Each root node
therefore denotes a valid summary. Apply the root-chain part of
Theorem~\ref{thm:pop-exact}.
\end{proof}

This certificate is naturally represented by shared Lean lemmas of type
$\operatorname{Pop}\ p\ A\ q$, not by a fully expanded list of transitions.
A final bridge can derive existential reachability from those lemmas without
normalizing an enormous concrete trace. The search chooses one derivation per
triple and slices away all nodes not needed by the selected root chain.

For debugging, the Python checker also computes the length of the represented
run:
\[
 L_i=1+\sum_{j\in\operatorname{children}(i)}L_j.
\]
Repeated child identifiers contribute repeatedly to the length, even though
their proofs are stored once. The local replay uses $O(K+\ell)$ node/reference
checks for $K$ stored nodes, in addition to input validation. The arithmetic
cost of the possibly large $L_i$ must not be hidden in a unit-cost claim. A Lean
checker need not compute these diagnostic lengths at all.

\begin{example}[An exponentially long witness in a linear-size DAG]
With one control, let
\[
 A_0\to\eps,\qquad A_i\to A_{i-1}A_{i-1}\quad(1\leq i\leq d).
\]
Starting from $A_d$, the only run to empty has length
$L_d=2^{d+1}-1$. There is one summary node per $A_i$; each nonbase node points
twice to its predecessor. At $d=30$, replay checks 31 nodes and reports
$2,147,483,647$ represented steps. The experiments expand the small cases
$d\leq10$ into concrete traces, but do not execute the billion-step run. The
large length is derived from the certificate recurrence.
\end{example}

\subsection{Negative certificate: a closed over-approximation}
A negative certificate supplies a finite relation
$C\subseteq Q\times\Gamma\times Q$. The checker independently enumerates the
consequences of every input rule and verifies that $C$ contains the pop bases
and is closed under both composition rules. It then runs the exact initial
stack query against $C$ and requires that no final control is admitted.

\begin{proposition}[Negative replay soundness]
If the exclusion checker accepts $C$, then no run in
\eqref{eq:pds-query} exists.
\end{proposition}
\begin{proof}
By induction on productive summary derivations, every actual summary belongs
to $C$. If a target run existed, its decomposition from
Theorem~\ref{thm:pop-exact} would supply a chain in $C$ ending in a final
control, contradicting the checker's last test.
\end{proof}

Here a larger relation is safe: it can prevent a negative certificate from
being accepted, but cannot create a false nonreachability proof. The polarity
is reversed for a positive claim. For example, the rule $A\to A$ does not pop
anything; nevertheless the relation containing $(p,A,p)$ is closed under that
rule. It would be an unsound positive witness, whereas the empty relation is a
valid exclusion for the pop query. A dedicated unit test submits precisely the
self-supporting positive node and verifies its rejection.

The negative checker does not call the agenda solver. Its deliberately simple
nested-loop audit costs
$O(r_0+r_1n+r_2n^2+\ell n^2+M)$ relation operations. It need not check that
$C$ is the \emph{smallest} closed relation. The solver emits the least relation
because that is enough to decide every query in the accepted fragment.

\begin{example}[Safety despite unbounded reachable stacks]
Let the initial control be $p$, with rules
$pA\to pAA$ and $pA\to p\eps$, and ask whether the empty stack can be reached
at a different control $q$. The reachable configurations include $pA^k$ for
arbitrarily large $k$. A finite enumeration of concrete configurations cannot
cover them all. The one-triple relation $C=\{(p,A,p)\}$ is nevertheless a closed
exclusion certificate for the target. A direct search with a height cap reports
truncation; the summary worker proves nonreachability. This is not a claim that
all finite abstractions fail---the relation $C$ is itself the useful finite
abstraction.
\end{example}

\subsection{Regular target languages, with an explicit bottom marker}
Empty-stack reachability is not the most convenient source interface. A user
may ask whether an error control can be reached with a stack matching a regular
pattern. The shipped compiler accepts a total DFA with states $D$, transition
$\delta:D\times\Gamma\to D$, accepting set $H$, and one initial DFA state
$i(q)$ for each source control. Its target predicate is
\[
  \operatorname{Target}(q,\beta)
  \quad\Longleftrightarrow\quad
  \delta^*(i(q),\beta)\in H.
\]
The DFA reads top to bottom. The original input's empty-stack final-control set
is deliberately ignored by this interface: the supplied regular target replaces
it, rather than being intersected with it silently.

The reduction appends a fresh bottom symbol $\bot$ to the initial stack and
adds fresh monitor controls $\bar d$ and a final control $\mathit{done}$. The
original rules cannot read or create $\bot$. For each original control $q$ and
symbol $A\in\Gamma\cup\{\bot\}$, add a one-way switch
\[
 (q,A)\to(\overline{i(q)},A).
\]
For $A\in\Gamma$, add
\[
 (\bar d,A)\to(\overline{\delta(d,A)},\eps),
\]
and for $d\in H$ add $(\bar d,\bot)\to(\mathit{done},\eps)$. There is no
return from a monitor control to an original control. The new target is empty
stack at $\mathit{done}$.

\begin{theorem}[Regular-target reduction]
The compiled system reaches its empty-stack target exactly when the original
system reaches a configuration satisfying $\operatorname{Target}$.
\end{theorem}
\begin{proof}
Given an original target run, append $\bot$ to its stacks, take the switch, and
consume the stack with the DFA monitor. Its final DFA state is accepting, so it
can remove $\bot$ and enter $\mathit{done}$.

Conversely, a compiled accepting run must enter the monitor, since only the
monitor can reach $\mathit{done}$. Before its first switch it uses only original
rules above the untouched bottom symbol. After that switch, its control is
determined by the DFA scan. The bottom pop can occur only after the original
stack has been consumed and the DFA has accepted it. The prefix before the
switch therefore corresponds to an original run reaching the regular target.
The bottom symbol also handles an already empty original stack, where no
ordinary source rule would be available.
\end{proof}

With $k=|D|$, the reduction has $n+k+1$ controls, $a+1$ stack symbols, and
$|\Delta|+n(a+1)+ka+|H|$ rules. The regression suite compares 100 compiled
queries with a complete concrete oracle on non-growing original systems. It
also distinguishes the words $01$ and $10$ to catch a reversed stack scan.
These tests are evidence about the compiler implementation, not a formal proof
of the reduction theorem in Lean.

\subsection{How Forge would use the worker}
An initial source adapter should accept a user-supplied pushdown model plus a
proved operational correspondence, rather than infer arbitrary program
semantics. Examples include a recursive interpreter with finitely many control
locations and stack-frame tags, a context-free parser with a finite alphabet,
or an error-state query for a recursive protocol. Once that adapter works,
recognized definitions can be translated automatically using reusable source
lemmas.

For a safety theorem, a forward simulation into the pushdown model can be
enough: every source error execution must map to an abstract error execution.
For a positive reachability theorem, an abstract witness must be concretizable;
a one-way over-approximation cannot justify it. Section~\ref{sec:bridges}
spells out the corresponding proof obligations rather than treating a shared
syntax tree as evidence that a translation is exact.
